The endpoint identification dichotomy is nevertheless complete. For each \(E\in\{A_Q\cap B_Q,B_Q\}\), Theorem \(\mathrm{thm:order\text{-}normal\text{-}factorization}\) proves the unique finite identity
\[
 \mathbf 1_E=\sum_{t:e_E(t)=1}\prod_{D\in\mathcal D(G)}\mathbf 1\{(R_v)_{v\in D}=t_D\},
\]
and every deleted order-inconsistent cell is null in every \(\mathfrak m\in\mathfrak M(G,M)\). Theorem \(\mathrm{thm:finite\text{-}oracle\text{-}complete\text{-}mctfid}\) proves
\[
 \operatorname{GID}_{G,M,\mathcal Z}(Q)=1
 \quad\Longleftrightarrow\quad
 \mathsf{MCTFID}(G,M,\mathcal Z,Q)=F\in\mathcal R(p),
\]
with \(F(P_{\mathcal Z}(\mathfrak m))=\theta_Q(\mathfrak m)\) for every compatible positive-denominator model. On the complementary branch, its exact certificate decodes to finite binary district-factorized signed-isotone models \(\mathfrak m_0,\mathfrak m_1\) satisfying
\[
 P_{\mathcal Z}(\mathfrak m_0)=P_{\mathcal Z}(\mathfrak m_1),\qquad
 P_{\mathfrak m_i}(B_Q)>0\ (i=0,1),\qquad
 \theta_Q(\mathfrak m_0)\ne\theta_Q(\mathfrak m_1).
\]
Thus finite identifying functionals and finite signed-isotone countermodels exist exactly on the two sides of the identification dichotomy; only the stronger district-local graphical compression and composable failure-trace theorem remain unresolved.